\section{Common Reporting Clock and Held-Out Episode Outcomes}
\label{app:heldout-recovery-protocol}
For source $c$ the bank supplies the onset and restoration scores
\[
 s_c^{\rm on}=\operatorname{logit}(M_c),\qquad
 s_c^{\rm off}=\operatorname{logit}(\Gamma_c-M_c).
\]
When the same source wins and exceeds threshold for six consecutive
monitored reports, the clock emits an alarm and names that source.
Six qualifying restoration scores for the named source close the episode.
The GLR and ExtraTrees supply their own scores to the same clock.

In the comparison with the GLR and ExtraTrees, each procedure completed
18 event and 24 normal test records, with 108 interventions in total.
Table~\ref{tab:closure-outcomes} gives the episode outcomes. A restoration
message can close an alarm that never matched a true intervention. Such
messages count separately from episodes still open when the record ends.

\begin{table}[t]\centering\small
\caption{Held-out event records at the selected thresholds. An orphan
restoration closes an unmatched episode. Open episodes are counted at
the end of the record.}
\label{tab:closure-outcomes}
\begin{tabular}{@{}lrrrr@{}}\toprule
Method & Timely & Extras & Orphan OFF & Open\\\midrule
History bank & \RecBankTimely & \RecBankExtras & \RecBankOrphanOff & \RecBankOpenEnd\\
Continuous GLR & \RecGlrTimely & \RecGlrExtras & \RecGlrOrphanOff & \RecGlrOpenEnd\\
ExtraTrees & \RecTreesTimely & \RecTreesExtras & \RecTreesOrphanOff & \RecTreesOpenEnd\\\bottomrule
\end{tabular}
\end{table}

All three procedures had zero alarms in \RecNormalHours{} monitored
normal hours. The one-sided 95\% Poisson rate bound was \RecNormalUpper{}
alarms/h for each.
